% =========================================================
\chapter{C06 --- Impulsions : le saut indiciel dans les circuits RC}
% =========================================================

\begin{rembox}[title=En bref]
\textbf{De quoi parle ce chapitre ?} Le couple $[R, C]$ revient partout (diodes,
transistors, AOP) : \textbf{volontairement}, pour réaliser des fonctions, ou
\textbf{involontairement}, comme modèle de phénomènes parasites qui limitent ou perturbent
le système. On l'étudie ici avec la deuxième catégorie de signaux : le \textbf{saut
indiciel} (les signaux carrés des horloges numériques). \\
\textbf{Les trois objectifs} :
\begin{enumerate}
  \item \textbf{le retour au temporel} : face à un saut brutal, le raccourci des complexes
        ne marche plus, on revient aux équations différentielles (premier ordre seulement :
        un seul condensateur, autant de résistances que nécessaire) ;
  \item \textbf{la continuité de la tension} : la tension aux bornes d'un condensateur ne
        peut pas sauter instantanément ($\Delta t\to0$) ; cela fixe les conditions
        initiales ;
  \item \textbf{la formule magique et sa recette} : une formule unique pour tous les circuits
        d'ordre 1 (passe-bas et passe-haut), qui dépend de la valeur initiale, de la valeur
        finale et de $\tau = RC$, pour résoudre les 4 \og angles\fg{} d'un signal carré.
\end{enumerate}
\end{rembox}

\section{Sauts indiciels (signaux carrés)}

\slide{c06-03}{Un créneau à l'entrée : ralentissement ou perturbation à la sortie.}

Les combinaisons R et C forment des filtres (volontaires ou non). Leur \textbf{analyse
temporelle} pour des sauts indiciels :
\begin{itemize}
  \item explique les \textbf{limites en fréquence des circuits numériques} (le créneau est
        \textbf{ralenti} : il monte et descend lentement) ;
  \item explique les \textbf{couplages parasites} qui produisent des \textbf{perturbations}
        (des pics à chaque front), par exemple quand on mélange numérique et analogique.
\end{itemize}
Il faut comprendre \textbf{une formule magique} et \textbf{une recette de cuisine}.

\section{Saut indiciel sur un passe-bas}

\slide{c06-04}{Avant et après la fermeture de l'interrupteur.}

